% Generated by scripts/assemble.py; corrections belong in source/corrections.json.
\examhead{135}{武汉理工大学 2024 年数学分析试卷}{129}
% p130-b01; source PDF page 130
\par\Needspace{5\baselineskip}
\blockanchor{p130-b01}
\majorquestion{一}计算题. 每题 10 分, 共 50 分.\par
% p130-b02; source PDF page 130
\blockanchor{p130-b02}
\question{1}计算极限 $\lim_{n\to \infty}(e^2 +6^n +8^n)^{\frac{1}{n}}$\par
% p130-b03; source PDF page 130
\blockanchor{p130-b03}
\question{2}设 $f(x) = (x^2 + a)e^x$ 在定义域 $\mathbb{R}$ 上有拐点但无极值点, 求 $a$ 的取值范围.\par
% p130-b04; source PDF page 130
\blockanchor{p130-b04}
\question{3}计算积分 $\int_0^{\frac{\pi}{2}}\frac{\cos x}{\sin x + \cos x}\mathrm{d}x.$\par
% p130-b05; source PDF page 130
\par\Needspace{6\baselineskip}
\blockanchor{p130-b05}
\question{4}计算曲面积分\par
\nopagebreak[4]
% p130-b06; source PDF page 130
\blockanchor{p130-b06}
\displayformula{\iint_ {\Sigma} x ^ {3} \mathrm{d} y \mathrm{d} z + y ^ {3} \mathrm{d} z \mathrm{d} x + (x ^ {2} + y ^ {2} + z ^ {3}) \mathrm{d} x \mathrm{d} y.}
% p130-b07; source PDF page 130
\blockanchor{p130-b07}
其中  $\Sigma: x^{2} + y^{2} + z^{2} = R^{2}$ ，方向取外侧.\par
% p130-b08; source PDF page 130
\blockanchor{p130-b08}
\question{5}讨论级数 $\sum_{n=1}^{\infty} \frac{x^n}{n(n+1)}$ 在 $[-1, 1]$ 上的一致收敛性, 并求其和函数.\par
% p130-b09; source PDF page 130
\par\Needspace{5\baselineskip}
\blockanchor{p130-b09}
\majorquestion{二}解答题. 每题 15 分, 共 60 分.\par
% p130-b10; source PDF page 130
\blockanchor{p130-b10}
\question{6}设 $f(x)$ 在 $x = 0$ 的某邻域内存在一阶连续导数, 且 $\lim_{x \to 0} \frac{f(x)}{x} = a > 0$. 讨论 $\sum_{n=1}^{\infty} (-1)^n f\left(\frac{1}{n}\right)$ 是否收敛, 若收敛是条件收敛还是绝对收敛?\par
% p130-b11; source PDF page 130
\blockanchor{p130-b11}
\question{7}已知 $z = f(x, y)$ 的全微分 $\mathrm{d}z = 2x\mathrm{d}x - 2y\mathrm{d}y$，且 $f(1, 1) = 2$.\par
% p130-b12; source PDF page 130
\blockanchor{p130-b12}
(1) 求 $f(x, y)$ 的表达式.\par
% p130-b13; source PDF page 130
\blockanchor{p130-b13}
(2) 求 $f(x, y)$ 在 $D: \left\{(x, y) \mid x^2 + \frac{y^2}{4} \leq 1\right\}$ 上的最大值与最小值.\par
% p130-b14; source PDF page 130
\par\Needspace{6\baselineskip}
\blockanchor{p130-b14}
\question{8}讨论函数\par
\nopagebreak[4]
% p130-b15; source PDF page 130
\blockanchor{p130-b15}
\displayformula{f (x, y) = \left\{ \begin{array}{l l} \frac {x y}{\sqrt {x ^ {2} + y ^ {2}}}, & (x, y) \neq (0, 0); \\ 0, & (x, y) = (0, 0). \end{array} \right.}
% p130-b16; source PDF page 130
\blockanchor{p130-b16}
在原点处的连续性与可微性.\par
% p130-b17; source PDF page 130
\par\Needspace{6\baselineskip}
\blockanchor{p130-b17}
\question{9}设 $f(x)$ 存在一阶连续导数, 且\par
\nopagebreak[4]
% p130-b18; source PDF page 130
\blockanchor{p130-b18}
\displayformula{f (x) = x + \int_ {0} ^ {x} t f ^ {\prime} (x - t) \mathrm{d} t + \int_ {0} ^ {+ \infty} e ^ {- \sqrt {x}} \mathrm{d} x.}
% p130-b19; source PDF page 130
\blockanchor{p130-b19}
求 $f(x)$ 的表达式.\par
% p130-b20; source PDF page 130
\par\Needspace{5\baselineskip}
\blockanchor{p130-b20}
\majorquestion{三}证明题. 每题10分, 共40分.\par
% p130-b21; source PDF page 130
\par\Needspace{6\baselineskip}
\blockanchor{p130-b21}
\question{10}证明如下问题:\par
\nopagebreak[4]
% p130-b22; source PDF page 130
\blockanchor{p130-b22}
\displayformula{\frac {x}{1 + x} <   \ln (1 + x) <   x (x > 0). \tag {1}}
% p130-b23; source PDF page 130
\blockanchor{p130-b23}
(2) $a_{n} = \sum_{k=1}^{n} \frac{1}{k} - \ln n$ 收敛.\par
% p130-b24; source PDF page 130
\blockanchor{p130-b24}
\question{11}用有限覆盖定理证明康拓定理, 即 $[a, b]$ 上的连续函数一致连续.\par
% p130-b25; source PDF page 130
\par\Needspace{6\baselineskip}
\blockanchor{p130-b25}
\question{12}设 $f(x)$ 在 $[-a, a]$ 上二阶连续可导，且 $f(0) = 0$，证明：存在 $\xi \in (-a, a)$，使得\par
\nopagebreak[4]
% p130-b26; source PDF page 130
\blockanchor{p130-b26}
\displayformula{f ^ {\prime \prime} (\xi) = \frac {1}{a ^ {2}} [ f (a) + f (- a) ].}
% p130-b27; source PDF page 130
\par\Needspace{6\baselineskip}
\blockanchor{p130-b27}
\question{13}已知连续函数列 $\{f_n(x)\}$ 在 $[a, b]$ 上一致收敛于 $f(x)$, 若 $\{x_n\} \subset [a, b]$ 满足 $\lim_{n \to \infty} x_n = c$, 证明:\par
\nopagebreak[4]
% p130-b28; source PDF page 130
\blockanchor{p130-b28}
\displayformula{\lim _ {n \rightarrow \infty} f _ {n} (x _ {n}) = f (c).}
